%!TEX root = ../../csuthesis_research.tex

\chapter{实验设计与预期创新}

本章围绕第四章给出的四类闭式解持续学习应用方法，
制定统一的实验设计与评测方案，
预先约定数据集、协议、对比方法、消融变量及结果呈现模板，
并对预期取得的效果给出说明。
需要指出的是，本章所述结果均为拟开展实验所对应的预期目标，
旨在为后续毕业论文阶段的实证研究提供可执行、可复现的工程基线。
本章内容与第三章的统一闭式解理论、
第四章的四类任务应用框架在记号、引用与流程上保持一致，
以确保整个研究脉络的连贯性。

\section{统一实验协议}

为保证不同任务、不同方法之间的可比性，拟在所有实验中采用一致的执行规范，具体包括以下几点。

\begin{itemize}
    \item \textbf{任务顺序与划分}：四类任务统一采用类别（或领域）增量协议；类别在每次实验前依据随机种子进行打乱，并按既定的类别数划分到若干任务中，确保同一方法在同一任务流上完成完整的增量过程；任务间类别互斥，避免出现新旧任务标签重叠造成的指标失真。
    \item \textbf{评测时机}：在每个增量任务结束后，立即对模型在所有已见任务的测试集上进行评测，得到任务级别的准确率矩阵或 mIoU 矩阵，用于平均准确率、遗忘率与后向迁移等指标的计算。
    \item \textbf{硬件与软件环境}：所有实验拟在一致的 GPU 计算环境（NVIDIA A100 80GB 或 V100 32GB）上完成，并使用相同的深度学习框架版本，以保证训练耗时、显存占用等效率指标具备可比性。
\end{itemize}

通过上述协议的统一约束，
期望实验结果能够更准确地反映方法本身的差异，
而非由实验设置带来的偶然偏差。
此外，所有训练日志、结果矩阵均拟以版本管理工具进行归档，
对应代码与配置文件统一开源，
以便后续研究者复现与扩展。

\section{数据集与任务设置}

为系统验证闭式解持续学习方法在不同任务上的通用性，
拟在四类任务上分别构造类别（或领域）增量基准，
覆盖二维医学影像、可穿戴时间序列、
二维与三维语义分割以及多模态领域终身学习。
各任务在数据规模、类别数量与模态复杂度上呈现明显差异，
有助于全面考察所提框架在不同应用场景下的泛化能力。

\subsection{医学图像疾病分类任务}

医学图像疾病分类任务拟选用 MedMNIST v2 数据集集合~\citep{yang2023medmnist}
中的三个有代表性的多类疾病分类数据集：
BloodMNIST、PathMNIST 与 OrganAMNIST。
该数据集在公开基准 LifeLonger~\citep{derakhshani2022lifelonger}
中已被广泛使用，并提供官方训练—验证—测试划分。
为构造类别增量协议，
拟把三个数据集分别划分为 4 个互斥任务，
使不同任务的类别两两不重叠；
图像大小统一调整为 $28\times28$，
输入模态为 RGB 二维灰度或彩色显微影像。
其中，BloodMNIST 覆盖外周血细胞显微图像，
PathMNIST 覆盖结肠癌组织病理切片，
OrganAMNIST 则覆盖腹部 CT 横切面器官分类，
三者在医学影像模态与诊断难度上具有明显的互补性。

\subsection{时间序列分类任务}

时间序列分类任务拟选用五个公开数据集：
UCI-HAR~\citep{ismi2016k}、UWave~\citep{liu2009uwave}、
DSA~\citep{altun2010human}、GRABMyo~\citep{pradhan2022multi}
与 WISDM~\citep{weiss2019wisdm}。
这五个数据集分别覆盖人体活动识别、手势识别、肌电信号识别等多种时序感知场景，
模态为多通道时间序列。
任务划分参照 \citet{qiao2024tscil} 的协议：
类别数较少的 UCI-HAR 与 UWave 按 3、4 个实验任务划分；
类别较多的 DSA、GRABMyo 与 WISDM 各划分若干实验任务用于评测，
并额外预留三个任务作为验证任务流以用于超参数搜索。
为评估方法对类别顺序的鲁棒性，
对每个数据集均预先生成五组不同的类别打乱顺序，
覆盖典型的“由易到难”与“由难到易”两种类别到达方式。

\subsection{语义分割任务}

语义分割任务覆盖二维图像与三维点云两类设置：

\begin{itemize}
    \item \textbf{二维图像}：Pascal VOC2012~\citep{everingham2010voc}，共 21 个语义类别（含背景），训练集 10\,582 张、验证集 1\,449 张。拟在 15-1（6 阶段）、15-5（2 阶段）与 10-1（11 阶段）三种增量协议下分别执行顺序、不相交与重叠三种设定下的实验。
    \item \textbf{三维点云}：S3DIS~\citep{armeni2016s3dis} 与 ScanNet~\citep{dai2017scannet}。S3DIS 包含 6 个区域共 272 个房间，Area 5 作为测试区域；ScanNet 包含 1\,513 个室内扫描场景。三维点云按滑动窗口分块，每块随机采样 2\,048 个点；类别顺序设置 $S^0$（按原始标注顺序）与 $S^1$（按字母顺序）两种。
\end{itemize}

\subsection{多模态领域终身学习任务}

多模态任务拟选用 MLLM-CL 基准~\citep{zhao2025mllm}，
任务流依次覆盖遥感、医疗、自动驾驶、科学与金融五个领域。
在落地层面，分别以 LLaVA-1.5~\citep{liu2023llava}
与 InternVL-Chat~\citep{chen2024internvl} 作为冻结主干，
并以每个领域对应一个 LoRA 专家的方式构造专家池；
路由头使用第三章给出的递归岭回归闭式解，
在领域顺序到达过程中完成增量更新。
为充分考察任务到达顺序的影响，
拟对五个领域的所有排列（120 种）做评测，
并报告对应指标的均值与标准差。
此外，路由头的输入特征拟从主干模型倒数第二层的隐藏状态中抽取，
以兼顾领域语义信息的丰富程度与计算开销之间的平衡。

\section{评价指标}

对应不同任务的特点，拟采用以下四类评价指标，分别衡量稳定性、可塑性、长期累积表现与工程效率。

\subsection{分类任务指标}

记 $a_{t,i}$ 为模型在完成第 $t$ 个任务训练后于第 $i$ 个任务上的测试准确率，则平均准确率 Acc 与平均遗忘率 Forg 分别定义为：

\begin{equation}
\label{eq:research_acc}
\mathrm{Acc}_T = \frac{1}{T}\sum_{i=1}^{T} a_{T,i},
\end{equation}

\begin{equation}
\label{eq:research_forg}
\mathrm{Forg}_T = \frac{1}{T-1}\sum_{i=1}^{T-1}\bigl(\max_{1\le t\le T-1} a_{t,i} - a_{T,i}\bigr).
\end{equation}

其中式~\eqref{eq:research_acc} 反映模型在所有已见任务上的整体可塑性，
式~\eqref{eq:research_forg} 度量旧任务被遗忘的程度。
除此之外，拟同时报告 Last-Acc（最后阶段全体类别的准确率）
与 Avg-Acc（任务推进过程中所有阶段准确率的均值），
以更细致地观察整体演化趋势。
对于类别不平衡的医学数据集，
拟额外报告宏平均准确率以避免少数类被多数类指标“淹没”。

\subsection{语义分割任务指标}

语义分割任务采用平均交并比（mIoU）作为基础指标，
并按照旧类、新类与全类三种范围分别报告：
旧类 mIoU 衡量知识保留能力，
新类 mIoU 衡量增量适应能力，
全类 mIoU 反映整体性能。
该指标与式~\eqref{eq:research_acc} 形式一致，
仅把任务级准确率 $a_{t,i}$ 替换为对应任务上的 mIoU。
对于点云任务，
另需单独报告每类 IoU 的方差，
以判别方法在不同类别上的稳健性。

\subsection{多模态路由任务指标}

多模态路由任务采用 MLLM-CL 基准约定的四个综合指标：末轮全任务平均性能 MFT、末轮新任务平均性能 MFN、平均累积准确率 MAA 以及后向迁移 BWT。其中 BWT 形式化为：

\begin{equation}
\label{eq:research_bwt}
\mathrm{BWT}_T = \frac{1}{T-1}\sum_{i=1}^{T-1}\bigl(a_{T,i} - a_{i,i}\bigr).
\end{equation}

BWT 越接近 0，
说明模型在学习新任务时对旧任务能力的破坏越小；
负值反映灾难性遗忘的存在。
除综合指标外，还需单独报告路由准确率，
以分析解析路由头在专家选择层面的鲁棒性。


\subsection{工程效率指标}

为反映方法的工程可行性，拟同时报告：
单任务增量更新时间、五任务累计训练总耗时、
显存峰值占用以及推理延迟。
各指标均在统一的硬件环境下测得，
并与可比基线在同一脚本框架下完成计时，
避免实现细节差异引入的偏差。


\section{对比方法}

为公平评估所提闭式解持续学习框架的优势，
拟选取四大类共十余种代表性方法作为基线，
覆盖回放、正则化、结构扩展与闭式解四个技术路线，
并针对语义分割与多模态路由两个细分任务额外选取专项基线，
具体如下。

\begin{itemize}
    \item \textbf{回放类方法}：iCaRL~\citep{rebuffi2017icarl}、End-to-End Incremental Learning（EEIL）~\citep{castro2018end}、Experience Replay（ER）~\citep{rolnick2019experience}、Dark Experience Replay（DER）~\citep{buzzega2020dark}、Herding~\citep{welling2009herding}。这类方法依赖一个历史样本缓冲区，是工业界最为常用的持续学习方案；将其作为基线可考察解析方法在无回放条件下能达到何种水平。
    \item \textbf{正则化类方法}：弹性权重巩固 EWC~\citep{kirkpatrick2017ewc}、记忆感知突触 MAS~\citep{aljundi2018memory}、知识蒸馏 LwF~\citep{li2017learning}、突触智能 SI~\citep{zenke2017continual}。这类方法通过对参数变化或输出分布施加正则项来约束遗忘，是与闭式解方法在“是否依赖梯度优化”这一维度上最直接对比的对象。
    \item \textbf{结构扩展类方法}：低秩适配 LoRA~\citep{hu2021lora}、硬注意力任务掩码 HAT~\citep{serra2018hat}、渐进式网络 Progressive Network~\citep{rusu2016progressive}、正交子空间 LoRA（O-LoRA）~\citep{wang2023orthogonal}、MoELoRA~\citep{chen2024coin}。这类方法通过为每个任务分配独立子结构来规避遗忘，可作为多模态路由任务的核心比较对象。
    \item \textbf{闭式解持续学习方法}：解析类增量学习 ACIL~\citep{zhuang2022acil}、双流解析 DS-AL~\citep{zhuang2024dsal}、解析持续学习 AIR~\citep{fang2024air}、生成性核线性扩展 GKEAL~\citep{GKEAL_Zhuang_CVPR2023}、广义解析类增量学习 GACL~\citep{GACL_Zhuang_NeurIPS2024}。这五条方法是直接同族对比基线，用于衡量本框架在统一性扩展与多任务适配方面的增益。
    \item \textbf{语义分割专项基线}：LwF-MC~\citep{rebuffi2017icarl}、ILT~\citep{ILT}、CIL~\citep{CIL}、MiB~\citep{cermelli2020mib}、PLOP~\citep{douillard2021plop}、SDR~\citep{SDR}、UCD+PLOP~\citep{UCD}、REMINDER~\citep{REMINDER}、RCIL~\citep{RCIL}、SPPA~\citep{SPPA}、CAF~\citep{CAF}、AWT+MiB~\citep{AWT}、EWF+MiB~\citep{EWF} 与 GSC~\citep{GSC}；三维点云另加入 \citet{Yang3D} 与 EWC~\citep{kirkpatrick2017ewc}，并对照综述~\citep{CSSsurvey} 中的结果。
    \item \textbf{多模态终身学习专项基线}：LoRA-FT~\citep{hu2021lora}、CL-MoE~\citep{huai2025cl}、HiDe~\citep{guo2025hide}、SEFE~\citep{chen2025sefe}、DISCO~\citep{guo2025federated} 等近期工作。各方法分别报告无回放与使用回放数据两个版本，以更全面地映射其对历史样本的依赖程度。
    \item \textbf{下界与上界参考}：所有任务均设置 Naive 基线（仅在当前任务上微调，不采取任何持续学习策略）作为下界，并以联合训练 Offline 作为理论上界，以便直观考察各方法相对的提升空间。
\end{itemize}

\section{消融实验设计}

为定量分析所提框架中各核心模块的贡献，
拟在四类任务上分别开展消融实验，
重点考察以下变量；
每条消融均严格在控制变量条件下完成，
其他超参数与协议保持与主结果实验一致。

\begin{itemize}
    \item \textbf{随机高维映射（RHL）}：以“是否引入随机隐藏层”作为开关变量，比较仅使用骨干编码器特征与引入 RHL 后的差异。预期在时间序列与语义分割任务上 RHL 显著改善特征的线性可分性，使新类别准确率有较大提升；对应实验汇总于第 5.6 节“消融表”中。
    \item \textbf{多尺度特征融合}：在时间序列与多模态主干上分别考察是否启用多尺度融合模块。预期融合后旧类别 mIoU 与平均准确率均有所提升，但代价是单任务计算量略有上升；该消融对应第 4.2 节时间序列框架的核心机制验证。
    \item \textbf{伪标签机制}：在二维与三维语义分割任务上，以“是否启用伪标签”作为开关变量。考虑到分割任务存在背景语义漂移，预期去掉伪标签后旧类 mIoU 显著下降；该实验同步在 Pascal VOC2012 与 S3DIS 上完成。
\end{itemize}


所得结论作为第四章框架设计选择的实证支撑。
为防止某一模块在不同任务上呈现相反的作用，


\section{预期结果与创新点}

基于第三章统一闭式解理论与第四章四类任务应用方法，
本节给出对各任务的预期结果与本课题的核心创新点。
预期结果以“对遗忘的抑制 + 准确率不低于代表性基线 + 增量阶段耗时显著降低”为统一目标，
并依据各任务的特点给出具体的量级估计；
预期创新点则从统一性、机制、任务扩展与工程四个维度展开。

\subsection{预期结果}

\begin{itemize}
    \item \textbf{医学图像疾病分类}：预期在各类医学影像数据集上，平均准确率不亚于主流回放类方法，并能在无回放的约束下显著优于主流正则化基线；平均遗忘率有望得到有效控制，从而在隐私敏感的医学场景中体现解析方法的实用价值。
    \item \textbf{时间序列分类}：预期在多类传感器数据集上，整体平均准确率能够贴近联合训练的理论上界，并进一步降低平均遗忘率；同时凭借随机高维映射与多尺度融合的协同作用，期望模型在面对类内变化复杂的数据分布时仍能保持良好的可塑性与稳定性。
    \item \textbf{语义分割}：预期在二维图像与三维点云的各类增量协议（如重叠、不相交设定）上，全类 mIoU 能够较现有代表性基线获得稳步提升，进一步逼近联合训练上界，充分证实闭式解在缓解稠密预测中语义漂移问题的优势。
    \item \textbf{多模态领域终身学习}：预期在冻结主干支持下，各综合评测评价指标稳居前列，后向迁移损失（遗忘现象）得到极大抑制，路由准确率维持在高水平，从而体现出闭式解路由相较于传统梯度路由在长序列稳定性上的优势。
    \item \textbf{工程效率}：得益于解析更新的设计，相比同等设置下的传统基线，预期增量阶段的总体训练耗时与显存占用均有显著下降，从而使整个框架更适配算力与存储受限的边缘下沉和部署场景。
\end{itemize}

\subsection{预期创新点}

\begin{itemize}
    \item \textbf{统一性创新}：将分类、分割与多模态专家路由三类原本独立处理的任务，纳入统一的递归岭回归闭式解主框架；其推导既兼容自相关矩阵—互相关矩阵的形式，也兼容递归更新的形式，为后续在其他模态上的扩展奠定理论基础。
    \item \textbf{机制创新}：以解析更新替代增量阶段的反向传播，从机制层面消除了由梯度冲突引发的灾难性遗忘；同时引入随机高维映射与多尺度融合，使冻结编码器架构在保持稳定性的前提下仍保留充分的可塑性。
    \item \textbf{任务扩展创新}：将闭式解持续学习从标准的二维分类问题成功扩展至语义分割（含背景语义漂移处理的伪标签机制）与多模态大模型专家路由（含跨领域 LoRA 专家组合），扩大了闭式解方法的应用边界。
    \item \textbf{工程创新}：构建了一套低开销、可复现、易部署的增量学习工程方案，仅需保存逆自相关矩阵 $\Psi_t$ 与解析分类头 $\widehat{\Phi}_t$ 即可完成持续更新，更符合医疗、工业与多模态边缘场景对存储、计算与隐私的综合约束。
\end{itemize}



\section{本章小结}

本章围绕第四章所述四类闭式解持续学习应用方法，
制定了从统一实验协议、数据集与任务设置、
评价指标、对比方法、消融实验设计
到结果输出模板的全流程实验方案，
并给出预期结果、核心创新点。
通过对评测协议的统一约束、对比方法的系统选取
以及消融变量的细致划分，
整个实验设计在工程可执行性、
统计严谨性与跨任务可对照性三方面均得到了保障。
下一章将基于本章给定的实验设计骨架，
进一步规划研究阶段的时间表、
阶段性里程碑，并对调研工作做整体总结与展望。

